
\documentclass[12pt, a4paper]{article}

% Cài đặt lề trang
\usepackage[a4paper, margin=2cm]{geometry}

% Cài đặt ngôn ngữ tiếng Việt và font chữ
\usepackage{fontspec}
\usepackage[vietnamese, bidi=basic, provide=*]{babel}
\babelfont{rm}{Noto Sans}

% Cài đặt bảng và màu sắc
\usepackage{tabularx}
\usepackage[table]{xcolor}

% Tắt lùi đầu dòng
\setlength{\parindent}{0pt}

% Chiều cao các hàng trong bảng
\renewcommand{\arraystretch}{2.2}

% Định nghĩa câu lệnh tạo khung cho 1 Task
\newcommand{\TaskTable}[1]{%
    \subsection*{Task #1: \underline{\hspace{10cm}}} % Tiêu đề Task có chừa chỗ viết tên Task
    \begin{tabularx}{\textwidth}{| c | p{3.2cm} | X | X |}
        \hline
        \rowcolor{gray!20}
        \textbf{STT} & \textbf{Thời gian} & \textbf{Nội dung} & \textbf{Công việc} \\
        \hline
        1 & & & \\[1.2cm]
        \hline
        2 & & & \\[1.2cm]
        \hline
        3 & & & \\[1.2cm]
        \hline
        4 & & & \\[1.2cm]
        \hline
    \end{tabularx}
    \vspace{0.8cm}
}

\begin{document}

\begin{center}
    \Large\textbf{KHUNG QUẢN LÝ 10 NHIỆM VỤ (TASKS)}
\end{center}

\vspace{0.5cm}

% ==================== TASK 1 ====================
\renewcommand{\arraystretch}{1.4} % Tối ưu giãn dòng cho bảng có chữ
\subsection*{Task 01: Khảo sát yêu cầu và thiết lập môi trường dự án}
\begin{tabularx}{\textwidth}{| c | p{3.2cm} | X | X |}
    \hline
    \rowcolor{gray!20}
    \textbf{STT} & \textbf{Thời gian} & \textbf{Nội dung} & \textbf{Công việc \& Phân công} \\
    \hline
    1 & 21/10/2026 - 23/10/2026 & Đọc hiểu tài liệu yêu cầu (Requirement Analysis) & Nghiên cứu đề bài Assignment 1, xác định phạm vi hệ thống và các chức năng chính. \newline \textbf{Phân công:} [Điền người thực hiện...] \\
    \hline
    2 & 24/10/2026 - 26/10/2026 & Khởi tạo dự án và quy chuẩn làm việc & Tạo Git repository, thống nhất luồng làm việc Git branching và quy ước commit. \newline \textbf{Phân công:} [Điền người thực hiện...] \\
    \hline
    3 & 27/10/2026 - 29/10/2026 & Họp phân công công việc & Phân chia vai trò, xác định rõ trách nhiệm của từng thành viên cho các Sprint tiếp theo. \newline \textbf{Phân công:} [Điền người thực hiện...] \\
    \hline

    4 &
    22/10/2026 &
    Thảo luận phân chia công việc bằng các issue&

    \textbf - Long thực hiện push các issue và setup CI cho các issue \\
    \hline
\end{tabularx}
\vspace{0.8cm}

% ==================== TASK 02 ====================
\subsection*{Task 02: API danh sách sản phẩm (GET /products)}

\begin{tabularx}{\textwidth}{|c|p{3.2cm}|Y|Y|}
    \hline
    \rowcolor{gray!20}
    \textbf{STT} &
    \textbf{Thời gian} &
    \textbf{Nội dung} &
    \textbf{Công việc \& Phân công} \\
    \hline

    1 &
    25/10/2026 &
    Xây dựng API \texttt{GET /products} &
    - Long thêm Prisma, model Product, migration và dữ liệu seed 6 sản phẩm, viết controller, service và test; tách PrismaService thành PrismaModule dùng chung \\
    \hline

    2 &
    25/10/2026 &
    Pull Request \#24 và CI &
    - Long mở Pull Request \#24 và merge vào main; CI trên main chạy thành công \\
    \hline

    3 &
    26/10/2026 &
    Bổ sung tài liệu API và kiểm thử &
    - Long thêm Swagger (\texttt{/docs}), collection Postman (4 request, 12 test) và ảnh minh họa cho 6 sản phẩm trong dữ liệu seed \\
    \hline

    4 &
    26/10/2026 &
    Kiểm tra hoàn thành Task 02 &
    - Long chạy lint (0 lỗi, 0 cảnh báo), test (3/3 qua), build và \texttt{docker compose up -{}-build}; \texttt{GET /products} trả về đủ 6 sản phẩm \\
    \hline
\end{tabularx}

\vspace{0.8cm}
\clearpage

% ==================== TASK 03 ====================
\subsection*{Task 03: Trang Menu (Next.js)}

\begin{tabularx}{\textwidth}{|c|p{3.2cm}|Y|Y|}
    \hline
    \rowcolor{gray!20}
    \textbf{STT} &
    \textbf{Thời gian} &
    \textbf{Nội dung} &
    \textbf{Công việc \& Phân công} \\
    \hline

    1 &
    24/10/2026 &
    Xây dựng trang menu &
    - Tài tạo trang menu, gọi API \texttt{GET /products}, hiển thị danh sách sản phẩm, có trạng thái loading và trạng thái empty \\
    \hline

    2 &
    25/10/2026 &
    Pull Request \#22 &
    - Tài mở Pull Request \#22 và merge vào main \\
    \hline

    3 &
    26/10/2026 &
    Sửa kiểu dữ liệu sản phẩm &
    - Tài sửa kiểu \texttt{Product} (trường \texttt{stock} không bắt buộc), Pull Request \#26 merge vào main, CI chạy thành công \\
    \hline

    4 &
    26/10/2026 &
    Kiểm tra hoàn thành Task 03 &
    - Long chạy \texttt{docker compose up -{}-build}, phát hiện cổng frontend map sai (4001 sang 3000 trong khi ứng dụng chạy ở 4001) và sửa lại; mở \texttt{localhost:4001} thấy menu 6 sản phẩm lấy từ API thật \\
    \hline
\end{tabularx}

\vspace{0.4cm}

\begin{tabularx}{\textwidth}{|Y|}
    \hline
    \rowcolor{gray!20}
    \textbf{Kết luận Sprint 1 (26/10/2026)} \\
    \hline
    Sprint 1 hoàn thành: Task 01, 02, 03 (issue \#3, \#4, \#5 đã đóng). Mục tiêu Sprint ``mở trình duyệt thấy menu quán, dữ liệu lấy từ API thật'' đã đạt, đã kiểm tra bằng lint, test, build và chạy Docker. \\
    \hline
\end{tabularx}

\vspace{0.8cm}

\TaskTable{04}
\clearpage

\TaskTable{05}
\TaskTable{06}
\clearpage

\TaskTable{07}
\TaskTable{08}
\clearpage

\TaskTable{09}
\TaskTable{10}

\end{document}
